\section{How it would be captured}
\label{sec:27.6}
Run the capture exercise of test 6 on these institutions, as on the charter.

\begin{itemize}
\item \emph{Defunding.} The Federal Art Project ended by appropriation. The answer is the permanent appropriation, ended only by a repeal on the record.
\item \emph{A content clause.} The 1990 clause is the precedent. The answer is that the endowment and the panels carry none.
\item \emph{The defense host.} The answer is the separation of \S\ref{sec:27.3}.
\item \emph{Funding by the studied.} The Royal Statistical Society's nominees were offered seats on Ofqual's advisory group only under a five-year agreement not to comment (\S\ref{sec:14.1}), and much of the discourse that examines these systems is paid for by the firms it examines (\S\ref{sec:15.6}). The answer is Appendix~\ref{sec:A}'s rule for the schools, extended: no institution in this chapter takes money from a party that builds, deploys or licenses models.
\item \emph{The revolving door.} The answer is a bar: no panelist has had a funding or contractual tie to such a party in the last three years, and the ties are published.
\item \emph{Access by grace.} A safe harbor a company grants is a company policy it can amend. The answer is the statute of \S\ref{sec:27.2}.
\item \emph{Staffing.} Institutions last because people staff them and a constituency stands behind them, and a party with an interest in the answer can build a pipeline of people as patiently as anyone (\S\ref{sec:18.1} traces one to a court's majority). The safeguards are the charter's: officers drawn by lot from member-nominated pools, caps on any one class or jurisdiction, and members free to leave.
\end{itemize}

The measure is the one \S\ref{sec:14.4} gives for capture without abolition: whether the cost of dissent is rising while the risk stays flat, and whether a panel's dissenters ever change what it decides.
